\chapter*{Abstract}
\addcontentsline{toc}{chapter}{Abstract}

\emph{The Zero Axis} asks a single structural question: if zero is not an
independent object but the vanishing of a relation, where can zero occur in a
normalized complementary relation?  The monograph adopts one foundational law.
Zero has no independent realization; zero is the vanishing of the canonical
defect of an admitted relation:
\[
\boxed{
\mathrm{A0}:\quad
R=0\Longrightarrow\mathfrak D_{\rm rel}[R]=0.
}
\]

Everything in the terminal proof layer is derived from A0 together with
definitions and standard analytic identities.

For a normalized binary contrast
\[
p(\sigma)=(\sigma,1-\sigma),\qquad 0<\sigma<1,
\]
exchange of the two sides induces the involution
\[
J(\sigma)=1-\sigma.
\]
Its unique fixed point is
\[
J(\sigma)=\sigma
\iff
\sigma=\frac12.
\]
Thus centered zero and the half are the same point in different coordinates:
\[
x:=\sigma-\frac12,\qquad
x=0\iff \sigma=\frac12.
\]

The second derivation is informational.  The Shannon defect
\[
D_H(\sigma)
=
\ln2-H_2(\sigma)
=
D_{\rm KL}\!\left(
(\sigma,1-\sigma)
\middle\|
(1/2,1/2)
\right)
\]
is non-negative and vanishes only at \(\sigma=1/2\).  Moreover
\[
D_H''(\sigma)=\frac1{\sigma(1-\sigma)}\ge4,
\]
so
\[
D_H(\sigma)\ge2\left(\sigma-\frac12\right)^2.
\]
Hence every transverse displacement from the relational zero carries a
strictly positive information defect.

The third derivation is projective.  With
\[
\Omega(s)=\frac{s}{1-s},
\qquad
\Delta_{\rm RC}(u)=|u^{-1}-\bar u|^2,
\]
one has the exact identity
\[
\boxed{
\Delta_{\rm RC}(\Omega(s))
=
\frac{4(\Re s-1/2)^2}{|s|^2|1-s|^2}
}
\qquad(s\ne0,1).
\]
Therefore the reciprocal--conjugation defect vanishes exactly on the Riemann
critical line.

Combining the Shannon bound with the projective identity gives the pointwise
coercive estimate
\[
\boxed{
D_H(\Re s)
\ge
\frac{|s|^2|1-s|^2}{2}\,
\Delta_{\rm RC}(\Omega(s)).
}
\]
The TIR \(U(1)\) phase/holonomy sector contributes additional non-negative
energy and therefore strengthens, rather than creates, this coercivity.

Applying A0 to a non-trivial completed-zeta zero \(\rho\) gives zero canonical
relational defect.  Hence
\[
0
\le
\frac{|\rho|^2|1-\rho|^2}{2}
\Delta_{\rm RC}(\Omega(\rho))
\le
E_{\rm rel}(\rho)
=
0,
\]
so
\[
\Delta_{\rm RC}(\Omega(\rho))=0
\quad\Longrightarrow\quad
\Re\rho=\frac12.
\]
Since \(\rho\) was arbitrary,
\[
\boxed{\mathrm{A0}\Longrightarrow\mathrm{RH}.}
\]

The proof is presented three times in complementary coordinates: relational
fixed-point geometry, Shannon/KL action cost, and projective/\(U(1)\)
reciprocal-conjugation geometry.  The formal Lean layer verifies the
pointwise-coercive endgame.  GREMLIN/OCTOPUS is recorded as the provenance
system used to compile dependency graphs, candidate proof routes, no-go scans,
and shortest-path reductions; public demonstrations are provided by the
GREMLIN-demo and CIEL--GREMLIN Benchmark repositories.

The older SOH-G and G024 chapters are retained because they record the
analytic route search, theorem-level no-go results, numerical certificates,
and falsification history that led to the minimum one-axiom proof graph.
